\ssscp{\tsc{Nuclear East Bomberai}}{11-05-02-01}
The \tsc{Nuclear East Bomberai} node contains \ili{Arguni}, \ili{Bedoanas}, and \ili{Erokwanas}.
Of all the \tsc{Tanimbar-Bomberai} languages,
the inclusion of these languages is the l\ili{east} certain.
Of the six instances of *j > \it{r} which define \tsc{Tanimbar-Bomberai},
possible reflexes have only been identified for three (of which
two are problematic), with a further four reflexes of *j identified in other words.
The putative reflexes of *j in \tsc{Nuclear East Bomberai} are
discussed below.

\begin{enumerate}
	\item	*ijuŋ could be linked with \ili{Arguni} \it{sigyó},
				\ili{Bedoanas}, \ili{Erokwanas} \it{suʔu} ‘nose’ with irregular loss of the
				initial vowel and irregular addition of a final vowel.
				This final vowel is unlikely to be epenthetic as the final epenthetic
				vowel is [e] in \ili{Arguni}
				and [a] in \ili{Bedoanas} and \ili{Erokwanas}.
				We would also need to posit irregular *u > \it{i} in \ili{Arguni} to
				derive \it{sigyó} from *ijuŋ.
				Thus, these words are unlikely to be connected.
				If they are, they show *j > \it{s}.
	\item	*pija > \ili{Arguni} \it{fís} ‘how much?’ could be a loan from
				\ili{Biak}, \ili{Kowiai}, or \ili{Geser} \it{fis}, or the final consonant may be a suffix
				(compare \ili{Onin} and \ili{Sekar} \it{firas}).
				If this is not a loan,
				it shows *j > \it{s} or *j > {\0}.
	\item	*qapəju > **ma-pəju > \ili{Arguni} \it{mari} is mostly unproblematic.
	\item	*pajay > \ili{Arguni} \it{pwásen},
				\ili{Bedoanas}, \ili{Erokwanas} \it{pása} ‘rice’ has irregular retention
				of *p = \it{p} where we expect *p > \it{f}.
				It is thus a loan,
				as it is in many other languages of the region.
	\item	*waRəj > \ili{Arguni},
				\ili{Bedoanas}, \ili{Erokwanas} \it{wárir} ‘rope’ is unproblematic.
	\item	*ŋajan > \ili{Arguni},
				\ili{Erokwanas} \it{gadin}, \ili{Bedoanas} \it{gadiŋ} ‘name’ is unproblematic.
	\item	*huaji > **ta-waji > \ili{Arguni} \it{twari} ‘friend’ is unproblematic.
\end{enumerate}

Of all these words,
only *qapəju > **ma-pəju > \ili{Arguni} \it{mari} ‘bitter’ has *j >
\it{r,} where other \tsc{Tanimbar-Bomberai} languages also have *j > **r.
The putative reflexes of *ijuŋ
and *pija with *j > \it{s} are similar to \tsc{Eastern Islands}
(\srf{11-03}) or SHWNG, but both have enough associated problems that they should probably
be set aside for the time being regarding subgrouping.
Thus, based on current evidence,
the basis for including these languages within \tsc{Tanimbar-Bomberai} is fragile.
Their position must be reassessed when more data is available.

Another possibility would be to place these languages within SHWNG.
However, the basis for that would be even more fragile.
The main defining change of SHWNG is merger of *j/*s > **s \citep[52]{Kamholz2014}
and there are only three putative cases of *j > \it{s} in \tsc{Nuclear
East Bomberai}, one of which is highly problematic,
and two of which are likely loans.
On the other hand,
there is one case of *j > \it{d}
and three cases of *j > \it{r},
none of which are problematic.

Further evidence that these languages do not belong within SHNWG
and do belong within STB comes from the merger of *d/*z > **d (\trf{11-02} and \trf{11-03}).
It may be that these languages should eventually be removed from
\tsc{Tanimbar-Bomberai} and placed as a primary branch of STB.
For the time being we include them within \tsc{Tanimbar-Bomberai}
as there is evidence for a link between these languages
and \tsc{\ili{Irarutu}-Kuri} (as discussed above),
and the evidence for including \tsc{\ili{Irarutu}-Kuri} within \tsc{Tanimbar-Bomberai}
is relatively good (see \srf{11-05-02-02}).

The evidence for placing \ili{Arguni},
\ili{Bedoanas}, and \ili{Erokwanas} within a single node comes mainly from
certain lexically specific changes in the vowels.
There are a handful of instances of final *a > \it{i},
two examples of *ə > \it{o},
and at l\ili{east} one example of final *i > \it{u}.
Examples are given in \trf{11-34}.

\begin{table}\tcp{\tsc{Nuclear East Bomberai} vowel changes}{11-34}
\begin{adjustbox}{width=\textwidth}
	\begin{threeparttable}\st{2pt}
	\begin{tabular}{llllllll}\lsptoprule
*gloss	&	name	&	moon	&	road, path	&	pig	&	paddle	&	skin	&	three\\
PMP	&	*ŋaj\tbr{a}n	&	*bul\tbr{a}n	&	*zal\tbr{a}n	&	*b\tbr{ə}Rək/*b\tbr{a}buy\su{†}	&	*b\tbr{ə}Rsay	&	*kul\tbr{i}t	&	*təlu\\	\midrule
\ili{Arguni}	&	{\hr}gád\tbr{i}n	&	{\hr}púr\tbr{i}n	&	{\hr}rár\tbr{i}n	&	{\hr}mb\tbr{o}	&	{\hr}pw\tbr{ó}res	&	{\hr}ár\tbr{u}t, ur\tbr{u}t	&	{\hr}táòr\\
\ili{Bedoanas}	&	{\hr}gad\tbr{i}ŋ	&	{\hr}úr\tbr{i}n	&	{\hr}rár\tbr{i}n	&	{\hr}mb\tbr{o}	&	{\hr}\tbr{ó}ras	&	{\hr}ár\tbr{u}t	&	{\hr}taur\\
Erokw.	&	{\hr}gad\tbr{i}n	&	{\hr}púr\tbr{i}n	&	{\hr}rór\tbr{i}n	&	{\hr}b\tbr{o}	&	{\hr}p\tbr{ó}ras	&	{\hr}ar\tbr{u}t	&	{\hr}tour\\
\ili{Irarutu}	&	\cg{(no)}	&	\cg{(seba)}	&	{\hr}rarn	&	\cg{(fne)}	&	{\hr}eʤe	&	{\hr}rit	&	{\hr}tor\\
\ili{Kuri}	&	\cg{(inɔi)}	&	\cg{(sarabɛ)}	&	{\hr}rarn	&	\cg{(bəne)}	&	{\hr}bɔyɛ	&	{\hr}riːtə̆	&	{\hr}tɔr\\
		\lspbottomrule
	\end{tabular}
		\begin{tablenotes}\tnsize
			\item[†] {These words for ‘pig’ could be from *bəRək or *babuy (see \trf{11-01}).
								The loss of the second consonant is unexplained under either hypothesis.
								Regardless of the source,
								this form for ‘pig’ is an innovation shared between these languages.}
		\end{tablenotes}
	\end{threeparttable}
\end{adjustbox}
\end{table}

There are also certain other lexically specific changes which
point to a shared ancestry for these three languages.
One example is apparent CV > VC metathesis in *təlu > \ili{Arguni}
\it{táòr}, \ili{Bedoanas} \it{taur}, and \ili{Erokwanas} \it{tour} ‘three’.

Additional evidence that these three languages have a shared
history comes from several items of shared vocabulary.
While many may be borrowings from neighbouring Papuan languages,
some have undergone regular sound changes in each language,
indicating that they were borrowed once in a common ancestor.
Three examples, with regular **g > \it{ʔ} in \ili{Bedoanas}
and \ili{Erokwanas} are: \ili{Arguni} \it{sigyó},
\ili{Bedoanas}, \ili{Erokwanas} \it{suʔu} ‘nose’, and \ili{Arguni} \it{sigye};
\ili{Bedoanas}, \ili{Erokwanas} \it{siʔi} ‘stomach’, as well as \ili{Arguni} \it{rigígit},
\ili{Bedoanas}, \ili{Erokwanas} \it{riríʔit} ‘white’.
